
\section{Introduction}
\label{sec:intro}

Sampling-based \ac{MPC} methods, such as \ac{CEM}--\ac{MPC}~\cite{pinneriSampleefficientCrossentropyMethod2021} and \ac{MPPI}~\cite{williamsInformationtheoretic2018,kazimRecentAdvancesPath2024}, have gained popularity for solving complex \acp{OCP} due to their ability to efficiently parallelize computations and handle arbitrary cost functions and system dynamics.

These methods iteratively refine a parameterized proposal \ac{PDF} over control sequences by drawing samples, evaluating their trajectory costs, and updating the proposal parameters accordingly.
Importance-weighting-based variants such as standard \ac{MPPI}~\cite{williamsInformationtheoretic2018}, and iterative extensions thereof~\cite{bhardwajSTORMIntegratedFramework2022} assign weights to the sampled control sequences based on their associated trajectory costs, thereby using a temperature parameter to control the sharpness of the weighting. 
This parameter is typically held constant or adjusted heuristically during optimization~\cite{pezzatoSamplingbasedModelPredictive2025}.
Similarly, \ac{CEM}-based methods rely on momentum-style smoothing~\cite{rubinsteinCrossEntropy2004} to prevent premature collapse of the proposal distribution.
These heuristics are often difficult to tune and may compromise convergence or performance across different problem settings.

Motivated by the success of trust region methods in \ac{RL}~\cite{schulmanTrustRegionPolicy2015,petersRelativeEntropyPolicy2010}, where principled \ac{KL}-constrained updates have largely replaced heuristic parameter tuning, we propose to transfer these ideas to sampling-based \ac{MPC}.
Specifically, for each proposal update, we constrain the new proposal to remain within a bounded divergence from the previous one. 
The resulting constrained optimization problem is solved via Lagrange multipliers, with optimal hyperparameters obtained from the dual formulation. 
The updated proposal is then projected back onto the original distribution family (\eg, Gaussian), yielding a feasible distribution from which control sequences are sampled in the subsequent iteration.

Orthogonal to the choice of update rule, the sampling method directly affects the quality and smoothness of the control.
In our previous work~\cite{walkerSampleefficientSmoothCrossentropy2026,arXiv26_Walker_IFAC}, we demonstrated that exchanging random samples with precomputed, optimally-placed deterministic samples based on \acp{LCD}~\cite{CDC09_HanebeckHuber} yield improved sample efficiency and smoother control sequences.
% This comes from their inherently homogeneous sample positioning, or more precisely, the smooth distribution of the size of the sample-free holes between them. 
% This property provides superior coverage and makes \ac{LCD}-based sampling particularly well-suited for integration and optimization, such as evaluating expectations in our trust region cost functions and constraining the updated proposal to stay close to the original.
\ac{LCD}-based sampling shows superior coverage and inherently positions the samples homogeneously, \ie, it induces a smooth distribution of the sizes of the sample-free gaps between them.
It is therefore well-suited for integration and optimization tasks.
In particular, it integrates well with trust region updates, which restrict the proposal to remain close to the original and thus to regions with high sample coverage.


% A trust region formulation naturally constrains the proposal update step, removing the need for manual temperature tuning or heuristic adaptation.

% \begin{figure}
%     \centering
%     % \includegraphics[width=0.95\columnwidth]{images/eye_catcher/gaussian_process_horizon.pdf} 
%     \caption{%
%     Eye-catcher.
%     }%
%     \label{fig:eye_catcher}
% \end{figure}
%
%
\paragraph*{Contribution}

Based on the success of trust region methods in \ac{RL}, we propose a trust region formulation for sampling-based \ac{MPC} that replaces heuristic proposal adaptations with principled \ac{KL}-constrained updates derived from dual optimization.
Building on our previous work~\cite{walkerSampleefficientSmoothCrossentropy2026,arXiv26_Walker_IFAC}, we further combine the trust region update with deterministic \ac{LCD}-based sampling to improve sample efficiency.
Both contributions are evaluated on two benchmark control tasks and compared against heuristic baselines across multiple sampling strategies, including random and low-discrepancy sampling (Sobol~\cite{joeConstructingSobolSequences2008} and Halton~\cite{owen2017randomizedhaltonalgorithmr}).

\paragraph*{Notation}
%
In this paper, underlined letters, \eg, $\vx$, denote vectors, boldface letters, such as $\rvx$, represent random variables, while boldface capital letters, such as $\mA$, indicate matrices.
% Sets are denoted by calligraphic letters, \eg, $\mathcal{E}$.
The expectation of a random variable is denoted by $\hat{\cdot}$, \eg, $\evx$, and covariance matrices are denoted by $\mC$.
% Diagonal matrices are represented by $\diag(\cdot)$.
% The indicator function is denoted by $\indFunc$, \eg, $\indFunc_{A}$ is \num{1} if event $A$ is true and \num{0} otherwise.
% Expectation operators are denoted by $\E{\cdot}$; \eg, $\E[p(\vx)]{\rvx}$ denotes the expectation of $\rvx$ under the \ac{PDF} $p(\vx)$. When appropriate, the \ac{PDF} argument is omitted for simplicity.
% The same convention applies to the variance operator, denoted by $\Var{\cdot}$.